\section*{Notation and Conventions}
\addcontentsline{toc}{section}{Notation and Conventions}
\label{sec:notation_and_conventions}

The following symbols are used consistently throughout this chapter.

\begin{table}[htbp]
    \centering
    \begin{tabular}{ll}
        \toprule
        Symbol & Meaning \\
        \midrule
        $t$              & Physical time \\
        $x(t)$           & Continuous state vector \\
        $z(t)$           & Algebraic variable vector \\
        $q(t)$           & Discrete mode or logic state \\
        $u(t)$           & Input vector \\
        $y(t)$           & Output vector \\
        $f(\cdot)$       & State-transition map \\
        $h(\cdot)$       & Output map \\
        $g(\cdot)$       & Algebraic constraint function \\
        $G(\cdot)$       & Implicit DAE residual \\
        $\gamma_j(\cdot)$& Event indicator (zero-crossing function) \\
        $t_e$            & Event time \\
        $(\cdot)^-$      & Value immediately before an event \\
        $(\cdot)^+$      & Value immediately after an event \\
        $(t,\nu)$        & Superdense time; $\nu \in \mathbb{N}_0$ is the micro-step index \\
        $i$              & Subsystem index \\
        $T_k$            & Communication point $k$ \\
        $H_k$            & Macro step size $H_k = T_{k+1} - T_k$ \\
        $\Delta t_{i,r}$ & Internal step size of subsystem $i$ at local step $r$ \\
        \bottomrule
    \end{tabular}
    \caption{Notation used throughout Chapter~\ref{chap:theoretical_background}.}
    \label{tab:notation}
\end{table}

In the physical and mathematical discussion, the term \emph{subsystem} denotes a constituent part of the modeled system.
In a tool-neutral co-simulation context, the executable software unit that encapsulates a subsystem model is called a \emph{simulation unit}.
In \emph{SysSimX}, the concrete realization of this concept is called a \emph{component} and is implemented through the \texttt{CoSimComponent} interface.
